% !TeX root = ../教材题库_全部条目_转写工作稿.tex
% Source: XJ-B1-C04-EX-P119-Y1481
% 第4章 / 4.3 对数函数 / 教材 p.119、120 / PDF 第126、127页
% OCR未校对：公式与文字均可能有误；未完成数学TeX转写；禁止直接发给学生
% 原题及相邻原解截图（核对用）：../题目原图/XJ-B1-C04-EX-P119-Y1481-01.webp, ../题目原图/XJ-B1-C04-EX-P119-Y1481-02.webp
% 本文件只在转写完成后作为题干使用；OCR 条目尚不是可靠题干。
\begin{bankitem}{4.3 例题 例题7}{第4章 / 4.3 对数函数 / 教材 p.119、120 / PDF 第126、127页}{XJ-B1-C04-EX-P119-Y1481}{OCR未校对：公式与文字均可能有误；未完成数学TeX转写；禁止直接发给学生}
\begin{ocrdraft}
[教材p.119 / PDF第126页]
\par
0二7]利用换底公式证明：
\par
（1）logab．lwa==1；
\par
（2）log\ocrglyph{201E}俨
\par
一log\ocrglyph{201E}b．
\par
77
\par
证明（1）由换底公式得，
\par
[教材p.120 / PDF第127页]
\par
loghb
\par
1
\par
logub==
\par
log,alog,a
\par
因此
\par
logab．loga==l.
\par
（2）由换底公式得，
\par
10力"一\_m
\par
由例7（2），有]力一
\par
\ocrglyph{00AF}log\ocrglyph{201E}b．
\par
logaannlog\ocrglyph{201E}a”
\par
俨，这个公式可以用
\par
来简化对数运算，如lc2
\end{ocrdraft}
\end{bankitem}
